% !TeX root = ./../../main.tex
% Miejsce ingerencji trzech rodzin metod w proces obliczeń (podrozdział 2.3, rys. fig:tb03:taxonomy).
% Rysunek własny (TikZ): górny wiersz to etapy obliczeń opartych na próbkowaniu,
% trzy pasy poniżej leżą pod etapami, które zmienia korekcja, mitygacja i redukcja błędów.
\begin{tikzpicture}[
  font=\footnotesize,
  stage/.style={draw, rounded corners=2pt, minimum height=0.8cm, text width=3.0cm, align=center, inner sep=2pt},
  band/.style={draw, fill=black!8, inner sep=2.5pt, align=center, minimum height=0.8cm},
  >={Stealth[length=2mm]}]
  % etapy obliczeń
  \node[stage] (circ) at (0,0)    {obwód kwantowy\\zadania};
  \node[stage] (meas) at (3.7,0)  {wykonania obwodu\\i~pomiar};
  \node[stage] (post) at (7.4,0)  {klasyczna obróbka\\statystyk};
  \node[stage] (res)  at (11.1,0) {wynik zadania};
  \draw[->] (circ) -- (meas);
  \draw[->] (meas) -- (post);
  \draw[->] (post) -- (res);
  % pasy: zakres ingerencji metod (szerokość pasa = etapy, które metoda zmienia)
  \node[band, text width=6.6cm] (qec) at (1.85,-1.05)
    {\textbf{korekcja (QEC):} nadmiarowe kubity,\\pomiar syndromu, operacje naprawcze};
  \node[band, text width=10.3cm] (qem) at (3.7,-2.0)
    {\textbf{mitygacja (QEM):} warianty obwodu lub obwody kalibracyjne,\\dodatkowe wykonania, obróbka zebranych statystyk};
  \node[band, text width=6.6cm] (red) at (9.25,-2.95)
    {\textbf{redukcja:} korekta wyniku modelem\\z~kalibracji; obwód zadania bez zmian};
\end{tikzpicture}
